\scanpage{103}
% archive sheet number: 102
\textcircled{3} \uline{Corollaire} $X$ schéma normal \deleted{\unsure{ou|et}} loc. intègre

\margin{$\begin{array}{ccccc}K&-&\widetilde K&-&K'\\ X&&\widetilde X&&X'\end{array}$}
$\unsure{D_i|D}$ diviseurs \inserted{de Cartier} \unsure{intègres} \deleted{\unsure{rédu}}

$X'$ revêtt \deleted{\unsure{galoisien}} normal galoisien, groupe $G$, tel que
les gpes d'inertie \inserted{$\deleted{\unsure{\mu_i}}I$ \unsure{(d'ordre $n_i$)}} relatifs \deleted{aux \unsure{comp.} sur} $D_i$ \unsure{soit|soient}
invariants \deleted{et commutent deux à deux, et que} \inserted{et que}
$X'$ soit modérément ramifié p.r. \deleted{aux \unsure{systèmes} des} \inserted{à} $D$. Donc on a \deleted{si $\xi$ est un pt géom. de $X$ et si on}
\[
1\mathbin{-}\overline I\mathbin{-}G\mathbin{-}\Gamma\mathbin{-}\deleted{\unsure{1}}\mathbin{-}1
\]
\deleted{$\overline I$ étant isomorphe au produit des $\mu(i)$.} De plus,
$X'/\overline I\deleted{\unsure{a}}=\widetilde X$ est \deleted{\unsure{donc}} un rev. \uline{principal} de $X$ de
groupe $\Gamma$. Donc la classe $\alpha\in H^2(\Gamma,I)$ de l'extension
définit un élément de $\alpha\in H^2(X,\widetilde X\overset{\Gamma}{\times} I)$, \deleted{D'ailleurs,}\inserted{et à fortiori
de $H^2(D,\widetilde D\overset{\Gamma}{\times}I)$. Or} on a un isomorphisme canonique
\[
\deleted{\unsure{I}}\widetilde D\overset{\Gamma}{\times}I\simeq\deleted{\prod\mu_{\underline n,D}}\mu_{\underline n,D}
\]
i.e. un isom. compatible avec les opérations de
$\Gamma$ \deleted{\unsure{de}} sur les faisceaux envisagés sur $(\widetilde X,\Gamma)$

or
\[
I_{\widetilde D}\ \deleted{\unsure{\simeq\mu_{n_i}}}=\mu_{\underline n,\widetilde D}
\]
\deleted{\unsure{Pour énoncer}} cet isomorphisme est tel que, en tant
que \unsure{fibré|fibre}, opérateurs $I_{\widetilde X}\simeq\mu_{\underline n}$, $X'$ est \deleted{\unsure{kummérien|kummerien}}
kummerien sur $\widetilde X$, relatif aux $\widetilde D_i$, : exposants $\underline n=(n_i)$.

Ceci dit, l'image de $\alpha$ dans $H^{\unsure{2|1}}(D,\mu_{\underline n})$ est égale à

\deleted{On peut supposer $X$ intègre. Le choix d'un
pt géométrique $\bar\xi_i$ de $X'$ au dessus un pt gén.
de $D_i$ définit un isom $I_i\simeq\mu$}

la classe induite sur $D$ par $c(\underline{\mathrm O}_X(D))\in H^2(X,\mu_{\underline n})$, i.e.
à $c(T_{D/X})\in H^2(D,\mu_{\underline n})$ \inserted{\uline{fibré normal}}
